% =====================================================================================
% CONCLUSIONS. Numbers from the result files quoted in results.tex and computationalMethodology.tex.
% 2026-10-29 (user): four paragraphs: (1) UrsaMLFF throughput, accuracy, MACE; (2) FA/MA tilt senses, r2SCAN vs PBEsol;
% (3) manufacturing and the band gap; (4) future work and the value for solar-cell manufacturing and design.
% Previous version: git history.
% =====================================================================================

{\changedcolor
We fine-tuned two classes of machine-learned potentials on PBEsol and \rSCAN{} labels of \CsSnI{}, \MASnI{} and
\FASnI{}. On held-out structures our in-house UrsaMLFF reaches errors of \SI{9.2e-5}{\hartree} per atom in the energy,
\SI{5.2e-4}{\hartree\per\bohr} in the forces and \SI{8.5e-7}{\hartree\per\bohr\cubed} in the stress. The fine-tuned
MACE models predict forces about twice as accurately but stresses 2.3 to 4.7 times less accurately, and they run 70 to
135 times slower: on one A100 GPU UrsaMLFF performs 960 to 2540 MD steps per second for cells of up to 2300 atoms, so
that the production MD of this work takes about \SI{10}{h} of GPU time instead of 40 to 59 days. With active-learning
labels for the molecular defects and close contacts that large-scale MD first reached, UrsaMLFF carries nanosecond MD
of thousands of atoms at meta-GGA accuracy, and the MACE models serve to validate it.

The \rSCAN{} potential reproduces the octahedral tilt patterns of the three compounds. Cooled from the cubic phase,
\CsSnI{} orders as $a^0a^0c^+$ and then $a^-a^-c^+$, \FASnI{} in-phase ($a^0a^0c^+$) and \MASnI{} antiphase
($a^0a^0c^-$), with single-tilt angles within \ang{0.3} of the deposited structures and tilt-ordering temperatures of
the hybrids within the measured ranges; it also forms the layered low-temperature phase of \FASnI{}, at
\SI{133}{K} against 125 to \SI{155}{K} measured. \changed{The low-temperature phase of \MASnI{} does not nucleate on
cooling, but a crystal prepared in it keeps it up to 160 to \SI{170}{K} on nanosecond time scales, against \SI{114}{K}
measured.} It describes the hybrids better than the PBEsol potential, which underestimates
the lattice constants by 2.0 to \SI{2.6}{\percent} (against 0.4 to \SI{1.0}{\percent}), leaves the antiphase tilt of
\MASnI{} weak and adds to \FASnI{} a second in-phase tilt that its single-crystal structure does not show; PBEsol places the transitions of
\CsSnI{} closer to experiment. Free-energy calculations at the experimental tilts explain the opposite senses: at
\SI{200}{K} \FASnI{} is lower in-phase by \changed{\SI{1.14 \pm 0.02}{\milli\hartree}} and \MASnI{} lower antiphase by
\changed{\SI{0.61 \pm 0.01}{\milli\hartree}} per formula unit. Antiphase tilting crosses the long axes of the lower and upper
faces of the A-site cage. FA, with an amine group at each end, spans the cage and can follow only one face, which costs
energy and orientational entropy; MA, with a single ammonium group, sits at one face and is lower in energy in the
crossed cage. Consistently, an in-phase \MASnI{} crystal converts to the antiphase tilt on heating, at \SI{190}{K}.
\changed{PBEsol gives the same signs, but its antiphase preference for \MASnI{} is weaker,
\SI{0.45}{\milli\hartree} per formula unit, too weak to order the tilt on cooling within the simulated time.}

Through the anneal and the cooling on the substrate, the manufacturing process leaves the films under an in-plane
tension of 50 to \SI{80}{MPa}, a strain of 0.2 to \SI{0.4}{\percent}, close to the levels measured in films, whether
the residual stress is taken from experiment or follows from cooling with the in-plane cell fixed at the anneal. The
constraint weakens the tilt of \CsSnI{} without changing its pattern. In the untilted hybrids the band gap opens with the expansion through
the elastic response alone, by 8 to \SI{15}{meV} with a deformation potential of \SI{3}{eV}; the changes of the thermal
ensembles, 0 to \SI{21}{meV}, agree with it but are not resolved against the scatter of the gap from frame to frame, 35
to \SI{67}{meV}. For \CsSnI{} the earlier PBEsol calculations give an opening of about \SI{20}{meV}.

Three extensions would bring the model closer to a device. Polycrystalline films contain grains and grain
boundaries, which relieve strain, nucleate tilt domains and may pin the tilt variant that forms on cooling; cells of
\num{20000} atoms already form such domains. An explicit substrate would replace the idealised boundary conditions by
the interface and the mismatch of thermal expansion along the actual cooling path. Spin--orbit coupling, strong for
iodine and tin, lowers the absolute gaps; the differences reported here cancel it to first order, but its effect on the
deformation potentials and on the tilt dependence of the gap remains to be tested. For manufacturers and designers of
tin perovskite solar cells, the approach offers a way to estimate, before a film is made, how the choice of cation, the
anneal temperature and the substrate shift the operating band gap and the mobility of iodide, and to compare
compositions and process windows on that basis.
}
